% =============================================================================
% COMP0496 --- Programação A
% Programação Orientada a Objetos --- Apresentação Completa
% Tema: "O Podrão --- Food Truck"
% Prof. Dr. André Yoshiaki Kashiwabara --- DCOMP/UFS
% =============================================================================

\documentclass[aspectratio=169]{beamer}
\usetheme{UFS}
\usepackage[utf8]{inputenc}
\usepackage[portuguese]{babel}
\usepackage{amsmath,amssymb}
\usepackage{graphicx}
\usepackage{booktabs}
\usepackage{listings}
\usepackage{xcolor}
\usepackage{tikz}
\usetikzlibrary{positioning,shapes.geometric}

% ---------------------------------------------------------------------------
% lstlisting style --- VS Code Light+ palette
% ---------------------------------------------------------------------------
\definecolor{bgcolor}{HTML}{FFFFFF}
\definecolor{linecolor}{HTML}{DDDDDD}
\definecolor{numcolor}{HTML}{999999}
\definecolor{kwcolor}{HTML}{0000FF}
\definecolor{strcolor}{HTML}{A31515}
\definecolor{cmtcolor}{HTML}{008000}
\definecolor{fncolor}{HTML}{795E26}
\definecolor{typecolor}{HTML}{267F99}

\lstdefinestyle{modernstyle}{
  language=Python,
  backgroundcolor=\color{bgcolor},
  rulecolor=\color{linecolor},
  numberstyle=\tiny\color{numcolor},
  keywordstyle=\color{kwcolor}\bfseries,
  stringstyle=\color{strcolor},
  commentstyle=\color{cmtcolor}\itshape,
  emph={print,len,range,super,isinstance,type,abs,
        receber_pagamento,ver_saldo,preco,total,__init__,__str__,__repr__,
        __eq__,__lt__,__hash__,__len__,__add__,__sub__,__mul__,simplificar,
        setter,deleter,from_csv_line,from_dict,calcular_taxa_servico,
        calcular_desconto,criar,registrar,adicionar,processar,finalizar,
        aplicar,gerar_relatorio,salvar,enviar_email},
  emphstyle=\color{fncolor},
  emph={[2]int,float,str,bool,list,dict,set,tuple,property,classmethod,
        staticmethod,ValueError,AttributeError,TypeError,None,True,False,
        Lanche,Pedido,Hamburguer,XTudo,XSalada,Bebida,
        DescontoPercentual,DescontoFixo,SemDesconto,LancheFactory,
        Cozinha,Caixa,RelatorioHTML,Armazenamento,Enviador,
        PedidoRuim,Retangulo,Quadrado,ABC,abstractmethod,
        dataclass,field},
  emphstyle={[2]\color{typecolor}},
  basicstyle=\ttfamily\scriptsize,
  frame=single,
  numbers=left,
  numbersep=8pt,
  tabsize=4,
  showstringspaces=false,
  breaklines=true,
  columns=fullflexible,
  literate={á}{{\'a}}1 {é}{{\'e}}1 {í}{{\'i}}1 {ó}{{\'o}}1 {ú}{{\'u}}1
           {ã}{{\~a}}1 {õ}{{\~o}}1 {ç}{{\c{c}}}1 {Ç}{{\c{C}}}1
           {ê}{{\^e}}1 {â}{{\^a}}1 {ô}{{\^o}}1 {à}{{\`a}}1,
}
\lstset{style=modernstyle}

% ---------------------------------------------------------------------------
% Metadata
% ---------------------------------------------------------------------------
\subtitle{COMP0496 -- Programação A \\ Programação Orientada a Objetos}
\title{Classes, Encapsulamento, Herança, Polimorfismo e Boas Práticas}
\author{Prof.\ Dr.\ André Yoshiaki Kashiwabara}

\begin{document}

% ===== TÍTULO =====
\begin{frame}[plain]
  \titlepage
\end{frame}

% ===== ROTEIRO =====
\begin{frame}[shrink=40]{Roteiro}
  \tableofcontents
\end{frame}

% #############################################################################
% PARTE 1: Classes, Objetos, Atributos e Métodos
% #############################################################################

% ############################################################
\section{Aquecimento --- Tudo em Python é Objeto}
% ############################################################

\begin{frame}[fragile]{Tudo em Python é Objeto}
  Você já usa objetos sem saber! Veja o que \texttt{type()} retorna:

  \begin{lstlisting}
>>> type("Oi")
<class 'str'>
>>> type([1, 2, 3])
<class 'list'>
>>> type(42)
<class 'int'>
  \end{lstlisting}

  \vspace{0.3cm}
  Cada valor em Python pertence a uma \textbf{classe} --- um tipo que define
  quais operações são possíveis.
\end{frame}

\begin{frame}[fragile]{Objetos Possuem Métodos}
  Métodos são funções que ``pertencem'' a um objeto:

  \begin{lstlisting}
>>> "programacao".upper()
'PROGRAMACAO'

>>> frutas = ["banana", "manga"]
>>> frutas.append("uva")
>>> frutas
['banana', 'manga', 'uva']
  \end{lstlisting}

  \vspace{0.3cm}
  \begin{block}{Ponto-chave}
    POO = criar \textbf{seus próprios tipos} (classes)
    com seus próprios métodos e atributos.
  \end{block}
\end{frame}

% ############################################################
\section{O Problema --- A Sopa de Variáveis}
% ############################################################

\begin{frame}[fragile]{O Cenário: Lanchonete ``O Podrão''}
  Você abriu um food truck e precisa de um sistema para gerenciar
  o cardápio.

  \vspace{0.3cm}
  Primeira tentativa: variáveis soltas.

  \begin{lstlisting}
lanche1_nome = "X-Tudo"
lanche1_preco = 18.50
lanche1_ingredientes = "pao, hamburguer, queijo, ovo, bacon"

lanche2_nome = "X-Salada"
lanche2_preco = 14.00
lanche2_ingredientes = "pao, hamburguer, queijo, alface, tomate"
  \end{lstlisting}
\end{frame}

\begin{frame}{Por que Isso é Ruim?}
  \begin{itemize}
    \item \textbf{Dados soltos}: nome, preço e ingredientes não estão
          conectados --- qualquer variável pode ser alterada isoladamente
    \item \textbf{Não escala}: com 50 lanches, seriam 150 variáveis!
    \item \textbf{Erros de digitação}: \texttt{lanche1\_preco} vs.\
          \texttt{lanche1\_prço} --- sem aviso
    \item \textbf{Sem comportamento}: não há como pedir a um lanche
          que ``se descreva''
  \end{itemize}

  \vspace{0.5cm}
  \begin{alertblock}{Analogia}
    É como anotar pedidos em guardanapos separados:
    fácil perder, fácil trocar, impossível de organizar.
  \end{alertblock}
\end{frame}

% ############################################################
\section{A Solução --- Classes e Objetos}
% ############################################################

\begin{frame}{Receita vs.\ Lanche Pronto}
  \begin{columns}[T]
    \begin{column}{0.48\textwidth}
      \textbf{Classe = Receita}
      \begin{itemize}
        \item Define \emph{quais} informações um lanche tem
        \item Define \emph{o que} um lanche pode fazer
        \item É um molde, um modelo
        \item Existe uma única vez no código
      \end{itemize}
    \end{column}
    \begin{column}{0.48\textwidth}
      \textbf{Objeto = Lanche Pronto}
      \begin{itemize}
        \item Preenche a receita com dados reais
        \item Cada objeto é independente
        \item Pode haver muitos objetos da mesma classe
        \item Vive na memória do programa
      \end{itemize}
    \end{column}
  \end{columns}
\end{frame}

\begin{frame}[fragile]{Primeira Classe: \texttt{Lanche}}
  \begin{lstlisting}
class Lanche:
    """Representa um lanche do cardapio."""

    def __init__(self, nome, preco, ingredientes):
        self.nome = nome
        self.preco = preco
        self.ingredientes = ingredientes

    def descrever(self):
        return f"{self.nome} - R${self.preco:.2f}"
  \end{lstlisting}
\end{frame}

\begin{frame}{Anatomia do \texttt{\_\_init\_\_}}
  \begin{itemize}
    \item \texttt{\_\_init\_\_} é o \textbf{construtor} --- é chamado
          automaticamente ao criar um objeto
    \item \texttt{self} é uma referência ao \textbf{próprio objeto}
          que está sendo criado
    \item \texttt{self.nome = nome} cria um \textbf{atributo de instância}
          --- cada objeto terá o seu
  \end{itemize}

  \vspace{0.3cm}
  \begin{block}{Atenção}
    \texttt{self} \textbf{sempre} é o primeiro parâmetro de qualquer método,
    mas \textbf{não} é passado na chamada --- Python cuida disso.
  \end{block}
\end{frame}

\begin{frame}[fragile]{Criando Objetos (Instâncias)}
  \begin{lstlisting}
xtudo = Lanche("X-Tudo", 18.50, "pao, hamburguer, queijo, ovo, bacon")
xsalada = Lanche("X-Salada", 14.00, "pao, hamburguer, queijo, alface")

print(xtudo.descrever())    # X-Tudo - R$18.50
print(xsalada.nome)         # X-Salada
print(xsalada.preco)        # 14.0
  \end{lstlisting}

  \vspace{0.3cm}
  Cada objeto é independente: alterar \texttt{xtudo.preco} não afeta
  \texttt{xsalada.preco}.
\end{frame}

\begin{frame}[fragile]{Comparação: Antes vs.\ Depois}
  \begin{columns}[T]
    \begin{column}{0.48\textwidth}
      \textbf{Sem OOP}
      \begin{lstlisting}[basicstyle=\ttfamily\scriptsize]
lanche1_nome = "X-Tudo"
lanche1_preco = 18.50
# ... e mais 148 variaveis
      \end{lstlisting}
    \end{column}
    \begin{column}{0.48\textwidth}
      \textbf{Com OOP}
      \begin{lstlisting}[basicstyle=\ttfamily\scriptsize]
xtudo = Lanche("X-Tudo", 18.50,
               "pao, hamburguer...")
# 1 variavel = 1 lanche completo
      \end{lstlisting}
    \end{column}
  \end{columns}

  \vspace{0.5cm}
  Dados e comportamento ficam \textbf{juntos} --- é isso que
  ``orientação a objetos'' significa.
\end{frame}

\begin{frame}{Vocabulário Essencial}
  \begin{table}
    \centering
    \begin{tabular}{ll}
      \toprule
      \textbf{Conceito} & \textbf{Significado} \\
      \midrule
      Classe         & Modelo/receita que define atributos e métodos \\
      Objeto         & Instância concreta de uma classe \\
      \texttt{\_\_init\_\_}  & Método construtor --- inicializa o objeto \\
      \texttt{self}  & Referência ao próprio objeto \\
      Atributo       & Variável que pertence ao objeto \\
      Método         & Função que pertence à classe \\
      \bottomrule
    \end{tabular}
  \end{table}
\end{frame}

\begin{frame}[fragile]{Métodos: Adicionando Comportamento}
  \begin{lstlisting}
class Lanche:
    def __init__(self, nome, preco, ingredientes):
        self.nome = nome
        self.preco = preco
        self.ingredientes = ingredientes

    def descrever(self):
        return f"{self.nome} - R${self.preco:.2f}"

    def aplicar_desconto(self, percentual):
        self.preco *= (1 - percentual / 100)

    def tem_ingrediente(self, ingrediente):
        return ingrediente in self.ingredientes
  \end{lstlisting}
\end{frame}

\begin{frame}[fragile]{Usando os Métodos}
  \begin{lstlisting}
xtudo = Lanche("X-Tudo", 18.50, "pao, hamburguer, queijo, ovo")

print(xtudo.descrever())             # X-Tudo - R$18.50

xtudo.aplicar_desconto(10)
print(xtudo.descrever())             # X-Tudo - R$16.65

print(xtudo.tem_ingrediente("ovo"))  # True
print(xtudo.tem_ingrediente("bacon"))# False
  \end{lstlisting}

  \vspace{0.3cm}
  O objeto \textbf{sabe} como se descrever, como aplicar desconto
  e como verificar ingredientes --- o comportamento está \emph{dentro} do objeto.
\end{frame}

% ############################################################
\section{Atributos de Classe vs Instância}
% ############################################################

\begin{frame}{Dois Tipos de Atributos}
  \begin{table}
    \centering
    \small
    \begin{tabular}{lp{4.5cm}p{4.5cm}}
      \toprule
      & \textbf{Atributo de Classe} & \textbf{Atributo de Instância} \\
      \midrule
      Onde define  & Dentro da classe, fora de métodos & Dentro de \texttt{\_\_init\_\_} \\
      Compartilhado & Sim, por todos os objetos & Não, cada objeto tem o seu \\
      Acesso       & \texttt{Classe.atributo} & \texttt{self.atributo} \\
      Exemplo      & Taxa de serviço (10\%) & Nome do lanche \\
      \bottomrule
    \end{tabular}
  \end{table}
\end{frame}

\begin{frame}[fragile]{Exemplo: Atributo de Classe}
  \begin{lstlisting}
class Lanche:
    taxa_servico = 0.10  # compartilhado por TODOS os lanches

    def __init__(self, nome, preco, ingredientes):
        self.nome = nome
        self.preco = preco
        self.ingredientes = ingredientes

    def preco_com_taxa(self):
        return self.preco * (1 + Lanche.taxa_servico)
  \end{lstlisting}

  \begin{lstlisting}
xtudo = Lanche("X-Tudo", 18.50, "pao, hamburguer, queijo")
print(xtudo.preco_com_taxa())  # 20.35

Lanche.taxa_servico = 0.15    # altera para TODOS
print(xtudo.preco_com_taxa())  # 21.275
  \end{lstlisting}
\end{frame}

\begin{frame}[fragile]{Armadilha: Atributo de Classe Mutável}
  \begin{lstlisting}
class Lanche:
    extras = []  # PERIGO! Lista compartilhada

    def __init__(self, nome):
        self.nome = nome

    def adicionar_extra(self, extra):
        self.extras.append(extra)
  \end{lstlisting}

  \begin{lstlisting}
a = Lanche("X-Tudo")
b = Lanche("X-Salada")
a.adicionar_extra("bacon")
print(b.extras)  # ['bacon'] -- b tambem tem bacon?!
  \end{lstlisting}
\end{frame}

\begin{frame}[fragile,shrink=10]{Correção: Criar Mutáveis no \texttt{\_\_init\_\_}}
  \begin{lstlisting}
class Lanche:
    def __init__(self, nome):
        self.nome = nome
        self.extras = []  # CORRETO: cada objeto tem sua lista

    def adicionar_extra(self, extra):
        self.extras.append(extra)
  \end{lstlisting}

  \begin{lstlisting}
a = Lanche("X-Tudo")
b = Lanche("X-Salada")
a.adicionar_extra("bacon")
print(a.extras)  # ['bacon']
print(b.extras)  # []  -- independente!
  \end{lstlisting}

  \begin{alertblock}{Regra de Ouro}
    Atributos mutáveis (listas, dicionários, conjuntos) devem ser
    criados em \texttt{\_\_init\_\_}, \textbf{nunca} como atributos de classe.
  \end{alertblock}
\end{frame}

% ############################################################
\section{Composição: Objetos Dentro de Objetos}
% ############################################################

\begin{frame}{Composição: Objetos Dentro de Objetos}
  Um \texttt{Cardapio} contém uma \textbf{lista de objetos}
  \texttt{Lanche}.

  \vspace{0.3cm}
  Isso é chamado de \textbf{composição}: um objeto é composto
  por outros objetos.

  \vspace{0.3cm}
  \begin{block}{Analogia}
    O cardápio do food truck não é um lanche --- ele \emph{contém}
    lanches. A relação é ``tem-um'' (has-a), não ``é-um'' (is-a).
  \end{block}
\end{frame}

\begin{frame}[fragile]{Classe \texttt{Cardapio}: Estrutura}
  \begin{lstlisting}
class Cardapio:
    """Gerencia a colecao de lanches do food truck."""

    def __init__(self, nome_food_truck):
        self.nome_food_truck = nome_food_truck
        self.lanches = []  # lista de objetos Lanche

    def adicionar(self, lanche):
        self.lanches.append(lanche)

    def mostrar(self):
        print(f"=== Cardapio: {self.nome_food_truck} ===")
        for lanche in self.lanches:
            print(f"  {lanche.descrever()}")
  \end{lstlisting}
\end{frame}

\begin{frame}[fragile]{Classe \texttt{Cardapio}: Métodos de Consulta}
  \begin{lstlisting}
class Cardapio:
    # ... (continuacao)

    def preco_medio(self):
        if not self.lanches:
            return 0
        total = sum(l.preco for l in self.lanches)
        return total / len(self.lanches)

    def mais_caro(self):
        if not self.lanches:
            return None
        return max(self.lanches, key=lambda l: l.preco)
  \end{lstlisting}
\end{frame}

\begin{frame}[fragile]{Usando o Cardápio}
  \begin{lstlisting}
cardapio = Cardapio("O Podrao")
cardapio.adicionar(Lanche("X-Tudo", 18.50, "pao, hamburguer..."))
cardapio.adicionar(Lanche("X-Salada", 14.00, "pao, alface..."))
cardapio.adicionar(Lanche("Misto Quente", 8.00, "pao, presunto"))

cardapio.mostrar()
  \end{lstlisting}

  Saída:
  \begin{lstlisting}[numbers=none]
=== Cardapio: O Podrao ===
  X-Tudo - R$18.50
  X-Salada - R$14.00
  Misto Quente - R$8.00
  \end{lstlisting}
\end{frame}

\begin{frame}[fragile]{Consultando o Cardápio}
  \begin{lstlisting}
print(f"Preco medio: R${cardapio.preco_medio():.2f}")
# Preco medio: R$13.50

campeao = cardapio.mais_caro()
print(f"Mais caro: {campeao.descrever()}")
# Mais caro: X-Tudo - R$18.50
  \end{lstlisting}

  \vspace{0.5cm}
  Note como \texttt{mais\_caro()} retorna um \textbf{objeto} \texttt{Lanche},
  e podemos chamar \texttt{descrever()} nele --- objetos conversam entre si!
\end{frame}

\begin{frame}{Diagrama: Composição}
  % Diagram: Composição Cardapio-Lanche
  % Coordinates:
  %   cardapio at (0, 0) -- objeto Cardapio
  %   lanche1 at (6, 1.5) -- primeiro Lanche
  %   lanche2 at (6, 0) -- segundo Lanche
  %   lanche3 at (6, -1.5) -- terceiro Lanche
  \begin{center}
    \begin{tikzpicture}[
      classbox/.style={draw, rounded corners, minimum width=3.5cm,
                       minimum height=1.2cm, text width=3.3cm, align=center,
                       fill=blue!10},
      objbox/.style={draw, rounded corners, minimum width=3.0cm,
                     minimum height=0.9cm, text width=2.8cm, align=center,
                     fill=green!10},
      >=stealth
    ]
      \node[classbox] (cardapio) at (0, 0)
        {\textbf{Cardapio}\\``O Podrão''};
      \node[objbox] (l1) at (7, 1.5)
        {\footnotesize X-Tudo\\R\$18.50};
      \node[objbox] (l2) at (7, 0)
        {\footnotesize X-Salada\\R\$14.00};
      \node[objbox] (l3) at (7, -1.5)
        {\footnotesize Misto Quente\\R\$8.00};

      \draw[->, thick] (cardapio) -- (l1)
        node[midway, above, font=\footnotesize] {contém};
      \draw[->, thick] (cardapio) -- (l2);
      \draw[->, thick] (cardapio) -- (l3);

      \node[font=\footnotesize, gray] at (0, -2.2)
        {1 objeto Cardapio};
      \node[font=\footnotesize, gray] at (7, -2.2)
        {N objetos Lanche};
    \end{tikzpicture}
  \end{center}
\end{frame}

% #############################################################################
% PARTE 2: Encapsulamento e Métodos Especiais
% #############################################################################

% ================================================================
\section{Encapsulamento}
% ================================================================

\begin{frame}[fragile]{O problema: estado inválido}
\begin{lstlisting}
xtudo = Lanche("X-Tudo", 18.50)
xtudo.preco = -999      # nenhum erro!
xtudo.preco = "gratis"  # nenhum erro!
print(xtudo.preco)      # "gratis"
\end{lstlisting}

\begin{alertblock}{Problema}
Qualquer código externo pode colocar o objeto em um estado que não faz sentido.
Um preço negativo ou textual viola as regras do nosso domínio.
\end{alertblock}

\medskip
Como proteger os dados internos de um objeto?
\end{frame}

\begin{frame}{Analogia: a caixa registradora}
\begin{columns}[T]
\column{0.55\textwidth}
Pense na caixa registradora do Podrão:
\begin{itemize}
    \item O \textbf{visor} mostra o total (acesso público)
    \item A \textbf{gaveta} só abre com a chave do operador (acesso restrito)
    \item Ninguém coloca a mão diretamente no dinheiro
\end{itemize}

\column{0.42\textwidth}
\begin{block}{Encapsulamento}
Esconder os detalhes internos de um objeto e expor apenas uma
interface controlada.
\end{block}
\end{columns}
\end{frame}

\begin{frame}{Convenções Python: público, protegido, privado}
Python \textbf{não tem} modificadores de acesso como Java.
Usa-se \textbf{convenções de nomenclatura}:

\medskip
\begin{center}
\begin{tabular}{lll}
\toprule
\textbf{Nome} & \textbf{Convenção} & \textbf{Significado} \\
\midrule
\texttt{nome}    & público   & Acesso livre \\
\texttt{\_nome}  & protegido & ``Use com cuidado'' \\
\texttt{\_\_nome}& privado   & Name mangling (dificulta acesso) \\
\bottomrule
\end{tabular}
\end{center}

\medskip
\begin{exampleblock}{Filosofia Python}
``We're all consenting adults here'' --- a linguagem confia no programador,
mas oferece mecanismos para sinalizar intenção.
\end{exampleblock}
\end{frame}

\begin{frame}[fragile]{Name mangling: como funciona}
Atributos com \texttt{\_\_} (duplo sublinhado) sofrem \textbf{name mangling}:

\begin{lstlisting}
class CaixaRegistradora:
    def __init__(self):
        self.__saldo = 0.0

    def receber_pagamento(self, valor):
        if valor > 0:
            self.__saldo += valor

    def ver_saldo(self):
        return self.__saldo
\end{lstlisting}
\end{frame}

\begin{frame}[fragile]{Name mangling: acesso externo}
\begin{lstlisting}
caixa = CaixaRegistradora()
caixa.receber_pagamento(25.00)
print(caixa.ver_saldo())   # 25.0

print(caixa.__saldo)
# AttributeError: 'CaixaRegistradora' object
#   has no attribute '__saldo'
\end{lstlisting}

\begin{alertblock}{Name mangling por dentro}
O Python renomeia \texttt{\_\_saldo} para \texttt{\_CaixaRegistradora\_\_saldo}.
\end{alertblock}

\begin{lstlisting}
# Funciona, mas NUNCA faca isso!
print(caixa._CaixaRegistradora__saldo)  # 25.0
\end{lstlisting}
\end{frame}

% ================================================================
\section{@property --- Acesso Controlado}
% ================================================================

\begin{frame}[fragile]{O problema dos getters e setters}
Em Java, usamos \texttt{getPreco()} e \texttt{setPreco()}.
Em Python, isso fica verboso e pouco natural:

\begin{lstlisting}
# Estilo Java em Python (evite!)
lanche.get_preco()
lanche.set_preco(20.0)

# Estilo Pythonico (queremos isto)
lanche.preco        # leitura
lanche.preco = 20.0 # escrita com validacao
\end{lstlisting}

A solução: \texttt{@property}.
\end{frame}

\begin{frame}[fragile]{\texttt{@property}: métodos disfarçados de atributos}
\begin{lstlisting}
class Lanche:
    def __init__(self, nome, preco):
        self._nome = nome
        self.preco = preco  # usa o setter!

    @property
    def preco(self):
        return self._preco

    @preco.setter
    def preco(self, valor):
        if valor < 0:
            raise ValueError("Preco negativo!")
        if valor > 500:
            raise ValueError("Preco acima do limite!")
        self._preco = valor
\end{lstlisting}
\end{frame}

\begin{frame}[fragile]{Usando a property}
\begin{lstlisting}
xtudo = Lanche("X-Tudo", 18.50)
print(xtudo.preco)   # 18.5 (chama o getter)

xtudo.preco = 22.0   # OK (chama o setter)
xtudo.preco = -10     # ValueError: Preco negativo!
xtudo.preco = 1000    # ValueError: Preco acima do limite!
\end{lstlisting}

\begin{block}{Vantagem}
O código externo usa sintaxe de atributo, mas internamente
há validação. Mudança transparente.
\end{block}
\end{frame}

\begin{frame}[fragile]{Property somente-leitura}
Use \texttt{@property} \textbf{sem} definir o setter:

\begin{lstlisting}
class Pedido:
    def __init__(self):
        self._itens = []

    def adicionar(self, lanche, qtd=1):
        self._itens.append((lanche, qtd))

    @property
    def total(self):
        return sum(l.preco * q for l, q in self._itens)
\end{lstlisting}

\begin{lstlisting}
pedido = Pedido()
pedido.adicionar(xtudo, 2)
print(pedido.total)     # 37.0
pedido.total = 0        # AttributeError!
\end{lstlisting}
\end{frame}

\begin{frame}[fragile]{Armadilha: recursão infinita}
\begin{alertblock}{Cuidado!}
Não use o mesmo nome para a property e o atributo interno.
\end{alertblock}

\begin{lstlisting}
class Errado:
    @property
    def preco(self):
        return self.preco  # chama a si mesmo!

    @preco.setter
    def preco(self, valor):
        self.preco = valor  # chama a si mesmo!
    # RecursionError!
\end{lstlisting}

\textbf{Solução:} armazene em \texttt{self.\_preco} (com prefixo).
\end{frame}

% ================================================================
\section{Métodos Especiais (Dunders)}
% ================================================================

\begin{frame}[shrink=10]{O que são dunders?}
Métodos com duplo sublinhado: \texttt{\_\_nome\_\_}

São chamados automaticamente pelo Python em situações específicas:

\medskip
\begin{center}
\begin{tabular}{ll}
\toprule
\textbf{Você escreve} & \textbf{Python chama} \\
\midrule
\texttt{print(obj)}      & \texttt{obj.\_\_str\_\_()} \\
\texttt{repr(obj)}       & \texttt{obj.\_\_repr\_\_()} \\
\texttt{obj1 == obj2}    & \texttt{obj1.\_\_eq\_\_(obj2)} \\
\texttt{obj1 < obj2}     & \texttt{obj1.\_\_lt\_\_(obj2)} \\
\texttt{len(obj)}        & \texttt{obj.\_\_len\_\_()} \\
\texttt{hash(obj)}       & \texttt{obj.\_\_hash\_\_()} \\
\texttt{obj1 + obj2}     & \texttt{obj1.\_\_add\_\_(obj2)} \\
\bottomrule
\end{tabular}
\end{center}
\end{frame}

\begin{frame}[fragile,shrink=10]{\texttt{\_\_str\_\_} vs \texttt{\_\_repr\_\_}}
\begin{lstlisting}
class Lanche:
    def __init__(self, nome, preco):
        self._nome = nome
        self._preco = preco

xtudo = Lanche("X-Tudo", 18.5)
print(xtudo)
# <__main__.Lanche object at 0x7f3b2c>  (feio!)
\end{lstlisting}

Agora com dunders:
\begin{lstlisting}
    def __str__(self):
        return f"{self._nome} - R${self._preco:.2f}"

    def __repr__(self):
        return f"Lanche('{self._nome}', {self._preco})"
\end{lstlisting}

\begin{lstlisting}
print(xtudo)        # X-Tudo - R$18.50  (amigavel)
repr(xtudo)         # Lanche('X-Tudo', 18.5) (tecnico)
\end{lstlisting}
\end{frame}

\begin{frame}[fragile]{\texttt{\_\_eq\_\_}: igualdade por valor}
Por padrão, dois objetos com os mesmos dados são \textbf{diferentes}:

\begin{lstlisting}
a = Lanche("X-Tudo", 18.5)
b = Lanche("X-Tudo", 18.5)
print(a == b)  # False (compara identidade!)
\end{lstlisting}

Solução:
\begin{lstlisting}
    def __eq__(self, other):
        if not isinstance(other, Lanche):
            return NotImplemented
        return (self._nome == other._nome
                and self._preco == other._preco)
\end{lstlisting}

\begin{lstlisting}
print(a == b)  # True
\end{lstlisting}
\end{frame}

\begin{frame}[fragile]{\texttt{\_\_lt\_\_}: ordenação com \texttt{sorted()}}
\begin{lstlisting}
    def __lt__(self, other):
        if not isinstance(other, Lanche):
            return NotImplemented
        return self._preco < other._preco
\end{lstlisting}

\begin{lstlisting}
cardapio = [
    Lanche("X-Tudo", 18.5),
    Lanche("Misto",   8.0),
    Lanche("Burgao", 22.0),
]
for l in sorted(cardapio):
    print(l)
# Misto - R$8.00
# X-Tudo - R$18.50
# Burgao - R$22.00
\end{lstlisting}
\end{frame}

\begin{frame}[fragile]{\texttt{\_\_hash\_\_}: objetos em conjuntos e dicionários}
Ao definir \texttt{\_\_eq\_\_}, o Python desabilita \texttt{\_\_hash\_\_}:

\begin{lstlisting}
lanches = {Lanche("X-Tudo", 18.5)}
# TypeError: unhashable type: 'Lanche'
\end{lstlisting}

Solução:
\begin{lstlisting}
    def __hash__(self):
        return hash((self._nome, self._preco))
\end{lstlisting}

\begin{alertblock}{Regra de ouro}
Se \texttt{a == b}, então \texttt{hash(a) == hash(b)}.
Use os mesmos atributos em ambos os métodos.
\end{alertblock}
\end{frame}

\begin{frame}[fragile]{\texttt{\_\_len\_\_}: tamanho de coleções}
\begin{lstlisting}
class Pedido:
    def __init__(self):
        self._itens = []

    def adicionar(self, lanche, qtd=1):
        self._itens.append((lanche, qtd))

    def __len__(self):
        return sum(q for _, q in self._itens)
\end{lstlisting}

\begin{lstlisting}
pedido = Pedido()
pedido.adicionar(Lanche("X-Tudo", 18.5), 2)
pedido.adicionar(Lanche("Misto", 8.0), 1)
print(len(pedido))  # 3
\end{lstlisting}
\end{frame}

\begin{frame}{Tabela resumo: principais dunders}
\begin{center}
\scriptsize
\begin{tabular}{lll}
\toprule
\textbf{Método} & \textbf{Operação} & \textbf{Quando usar} \\
\midrule
\texttt{\_\_init\_\_}  & Construtor          & Sempre \\
\texttt{\_\_str\_\_}   & \texttt{print()}, \texttt{str()} & Saída amigável \\
\texttt{\_\_repr\_\_}  & \texttt{repr()}, terminal & Debug, logs \\
\texttt{\_\_eq\_\_}    & \texttt{==}         & Comparação por valor \\
\texttt{\_\_lt\_\_}    & \texttt{<}, \texttt{sorted()} & Ordenação \\
\texttt{\_\_hash\_\_}  & \texttt{set}, \texttt{dict} & Após definir \texttt{\_\_eq\_\_} \\
\texttt{\_\_len\_\_}   & \texttt{len()}      & Coleções \\
\texttt{\_\_add\_\_}   & \texttt{+}          & Soma/concatenação \\
\texttt{\_\_sub\_\_}   & \texttt{-}          & Subtração \\
\texttt{\_\_mul\_\_}   & \texttt{*}          & Multiplicação \\
\bottomrule
\end{tabular}
\end{center}
\end{frame}

% ================================================================
\section{Prática --- Classe Fracao}
% ================================================================

\begin{frame}{Classe \texttt{Fracao}: especificação}
Vamos construir uma classe completa que usa \textbf{tudo} que aprendemos:

\begin{itemize}
    \item Validação no construtor (denominador $\neq 0$)
    \item Simplificação automática via MDC
    \item Properties somente-leitura
    \item \texttt{\_\_str\_\_}, \texttt{\_\_repr\_\_}
    \item \texttt{\_\_eq\_\_}, \texttt{\_\_lt\_\_}, \texttt{\_\_hash\_\_}
    \item \texttt{\_\_add\_\_}, \texttt{\_\_sub\_\_}, \texttt{\_\_mul\_\_}
\end{itemize}
\end{frame}

\begin{frame}[fragile,shrink=15]{Classe \texttt{Fracao}: construtor e simplificação}
\begin{lstlisting}
from math import gcd

class Fracao:
    def __init__(self, num, den=1):
        if den == 0:
            raise ValueError("Denominador nao pode ser zero!")
        # sinal sempre no numerador
        if den < 0:
            num, den = -num, -den
        divisor = gcd(abs(num), den)
        self._num = num // divisor
        self._den = den // divisor

    @property
    def num(self):
        return self._num

    @property
    def den(self):
        return self._den
\end{lstlisting}
\end{frame}

\begin{frame}[fragile]{Classe \texttt{Fracao}: representação}
\begin{lstlisting}
    def __str__(self):
        if self._den == 1:
            return str(self._num)
        return f"{self._num}/{self._den}"

    def __repr__(self):
        return f"Fracao({self._num}, {self._den})"
\end{lstlisting}

\begin{lstlisting}
print(Fracao(4, 6))   # 2/3
print(Fracao(6, 3))   # 2
repr(Fracao(4, 6))    # Fracao(2, 3)
\end{lstlisting}
\end{frame}

\begin{frame}[fragile,shrink=20]{Classe \texttt{Fracao}: comparação e hash}
\begin{lstlisting}
    def __eq__(self, other):
        if not isinstance(other, Fracao):
            return NotImplemented
        return (self._num == other._num
                and self._den == other._den)

    def __lt__(self, other):
        if not isinstance(other, Fracao):
            return NotImplemented
        return self._num * other._den < other._num * self._den

    def __hash__(self):
        return hash((self._num, self._den))
\end{lstlisting}
\end{frame}

\begin{frame}[fragile,shrink=15]{Classe \texttt{Fracao}: operações aritméticas}
\begin{lstlisting}
    def __add__(self, other):
        if not isinstance(other, Fracao):
            return NotImplemented
        return Fracao(self._num * other._den
                      + other._num * self._den,
                      self._den * other._den)

    def __sub__(self, other):
        if not isinstance(other, Fracao):
            return NotImplemented
        return Fracao(self._num * other._den
                      - other._num * self._den,
                      self._den * other._den)

    def __mul__(self, other):
        if not isinstance(other, Fracao):
            return NotImplemented
        return Fracao(self._num * other._num,
                      self._den * other._den)
\end{lstlisting}
\end{frame}

\begin{frame}[fragile]{Classe \texttt{Fracao}: demonstração}
\begin{lstlisting}
a = Fracao(1, 2)
b = Fracao(1, 3)

print(a + b)          # 5/6
print(a - b)          # 1/6
print(a * b)          # 1/6
print(a == Fracao(2, 4))  # True (simplificado!)
\end{lstlisting}

\begin{lstlisting}
fracs = [Fracao(3,4), Fracao(1,2), Fracao(2,3)]
print(sorted(fracs))  # [1/2, 2/3, 3/4]
\end{lstlisting}

\begin{lstlisting}
# set() funciona porque temos __eq__ e __hash__
s = {Fracao(1,2), Fracao(2,4), Fracao(3,6)}
print(s)              # {1/2}
\end{lstlisting}
\end{frame}

% #############################################################################
% PARTE 3: Herança, Polimorfismo e Composição
% #############################################################################

% =============================================================================
\section{Herança Simples}
% =============================================================================

\begin{frame}[fragile]{O Problema: Código Duplicado}
  No food truck \textbf{O Podrão}, temos três lanches com código repetido:

  \begin{columns}[T]
    \column{0.32\textwidth}
    \begin{lstlisting}
class XTudo:
    def __init__(self, preco):
        self.preco = preco
        self.nome = "X-Tudo"

    def descrever(self):
        return f"{self.nome}: R${self.preco:.2f}"

    def preparar(self):
        print("Fritando hambúrguer...")
    \end{lstlisting}

    \column{0.32\textwidth}
    \begin{lstlisting}
class Dogao:
    def __init__(self, preco):
        self.preco = preco
        self.nome = "Dogão"

    def descrever(self):
        return f"{self.nome}: R${self.preco:.2f}"

    def preparar(self):
        print("Montando hot-dog...")
    \end{lstlisting}

    \column{0.32\textwidth}
    \begin{lstlisting}
class Acai:
    def __init__(self, preco):
        self.preco = preco
        self.nome = "Açaí"

    def descrever(self):
        return f"{self.nome}: R${self.preco:.2f}"

    def preparar(self):
        print("Batendo açaí...")
    \end{lstlisting}
  \end{columns}
\end{frame}

\begin{frame}{Por que Código Duplicado é Ruim?}
  \begin{alertblock}{Violação do Princípio DRY}
    \textbf{Don't Repeat Yourself} --- código duplicado é fonte de bugs.
  \end{alertblock}

  \bigskip

  \begin{itemize}
    \item \texttt{\_\_init\_\_} e \texttt{descrever} são \textbf{idênticos} nas três classes
    \item Se precisar adicionar \texttt{calorias}: alterar \textbf{três} classes
    \item Se corrigir um bug em \texttt{descrever}: corrigir em \textbf{três} lugares
    \item E se o cardápio crescer para 15 itens?
  \end{itemize}

  \bigskip

  \begin{exampleblock}{Solução}
    Extrair o que é \textbf{comum} para uma classe base $\Rightarrow$ \textbf{Herança}.
  \end{exampleblock}
\end{frame}

\begin{frame}{O Relacionamento ``É Um''}
  Herança modela o relacionamento \textbf{``é um''} (\emph{is-a}):

  \medskip

  \resizebox{\textwidth}{!}{%
  \begin{tabular}{lcc}
      \toprule
      \textbf{Afirmação} & \textbf{Verdade?} & \textbf{Herança?} \\
      \midrule
      XTudo é um Lanche           & \checkmark & Sim \\
      Dogão é um Lanche           & \checkmark & Sim \\
      Açaí é um Lanche            & \checkmark & Sim \\
      Pedido é uma CaixaRegistradora & $\times$ & Não \\
      Pedido \emph{tem} Lanches   & \checkmark & Composição \\
      \bottomrule
    \end{tabular}%
  }% end resizebox

  \medskip

  \begin{alertblock}{Regra de Ouro}
    Se ``B é um A'' faz sentido no domínio $\Rightarrow$ \texttt{class B(A)}.
  \end{alertblock}
\end{frame}

\begin{frame}{Hierarquia de Classes}
  % Coordinate map:
  % (0,3)    = Lanche (parent node)
  % (-3,0)   = XTudo (child left)
  % (0,0)    = Dogao (child center)
  % (3,0)    = Acai (child right)
  \begin{center}
    \begin{tikzpicture}[scale=1.1, every node/.style={scale=1.1}]
      \node[draw, fill=blue!15, rounded corners, minimum width=3cm, minimum height=1cm, font=\bfseries]
        (lanche) at (0,3) {Lanche};
      \node[draw, fill=green!10, rounded corners, minimum width=2.5cm, minimum height=1cm]
        (xtudo) at (-4,0) {XTudo};
      \node[draw, fill=green!10, rounded corners, minimum width=2.5cm, minimum height=1cm]
        (dogao) at (0,0) {Dogão};
      \node[draw, fill=green!10, rounded corners, minimum width=2.5cm, minimum height=1cm]
        (acai) at (4,0) {Açaí};

      \draw[->, thick] (xtudo) -- node[left, font=\scriptsize]{herda} (lanche);
      \draw[->, thick] (dogao) -- node[right, font=\scriptsize]{herda} (lanche);
      \draw[->, thick] (acai) -- node[right, font=\scriptsize]{herda} (lanche);

      \node[below=0.2cm of xtudo, font=\scriptsize\ttfamily, text=gray] {class XTudo(Lanche)};
      \node[below=0.2cm of dogao, font=\scriptsize\ttfamily, text=gray] {class Dogao(Lanche)};
      \node[below=0.2cm of acai, font=\scriptsize\ttfamily, text=gray] {class Acai(Lanche)};
    \end{tikzpicture}
  \end{center}
\end{frame}

\begin{frame}[fragile]{Classe Base: \texttt{Lanche}}
  \begin{lstlisting}
class Lanche:
    def __init__(self, nome, preco):
        self.nome = nome
        self.preco = preco

    def descrever(self):
        return f"{self.nome}: R${self.preco:.2f}"

    def preparar(self):
        print(f"Preparando {self.nome}...")
  \end{lstlisting}

  \begin{exampleblock}{O que mudou?}
    Todo código comum agora vive em \textbf{um único lugar}.
  \end{exampleblock}
\end{frame}

\begin{frame}[fragile]{Subclasses com \texttt{super()}}
  \begin{lstlisting}
class XTudo(Lanche):
    def __init__(self, preco, molho="especial"):
        super().__init__("X-Tudo", preco)
        self.molho = molho

    def preparar(self):
        print(f"Fritando hambúrguer com molho {self.molho}...")
  \end{lstlisting}

  \begin{lstlisting}
class Dogao(Lanche):
    def __init__(self, preco, tamanho="grande"):
        super().__init__("Dogão", preco)
        self.tamanho = tamanho

    def preparar(self):
        print(f"Montando hot-dog {self.tamanho}...")
  \end{lstlisting}
\end{frame}

\begin{frame}[fragile]{Herdando Métodos Automaticamente}
  \begin{lstlisting}
x = XTudo(25.0)
d = Dogao(18.0)

# descrever() é herdado de Lanche --- funciona sem reescrever!
print(x.descrever())   # X-Tudo: R$25.00
print(d.descrever())   # Dogão: R$18.00

# preparar() foi sobrescrito em cada subclasse
x.preparar()   # Fritando hambúrguer com molho especial...
d.preparar()   # Montando hot-dog grande...
  \end{lstlisting}

  \begin{exampleblock}{Regra}
    Se a subclasse \textbf{não} redefine um método, usa o da superclasse.
  \end{exampleblock}
\end{frame}

\begin{frame}[fragile]{Verificando Tipos com \texttt{isinstance()}}
  \begin{lstlisting}
x = XTudo(25.0)

isinstance(x, XTudo)    # True
isinstance(x, Lanche)   # True  (XTudo EH UM Lanche)
isinstance(x, Dogao)    # False

type(x) == XTudo         # True
type(x) == Lanche        # False (type() é exato)
  \end{lstlisting}

  \bigskip

  \begin{center}
    \begin{tabular}{lcc}
      \toprule
      \textbf{Verificação} & \textbf{\texttt{isinstance}} & \textbf{\texttt{type() ==}} \\
      \midrule
      Considera herança? & Sim & Não \\
      Uso recomendado    & Geral & Raro \\
      \bottomrule
    \end{tabular}
  \end{center}
\end{frame}

% =============================================================================
\section{Sobrescrita de Métodos}
% =============================================================================

\begin{frame}[fragile]{Override: Redefinindo um Método}
  \textbf{Sobrescrita} (\emph{override}): a subclasse fornece sua própria versão de um método herdado.

  \begin{lstlisting}
class Acai(Lanche):
    def __init__(self, preco, complementos=None):
        super().__init__("Açaí", preco)
        self.complementos = complementos or ["granola", "banana"]

    def preparar(self):
        # Sobrescreve completamente o método de Lanche
        print(f"Batendo açaí com {', '.join(self.complementos)}...")
  \end{lstlisting}

  \begin{alertblock}{Atenção}
    O método da subclasse \textbf{substitui} o da superclasse --- não são acumulados.
  \end{alertblock}
\end{frame}

\begin{frame}[fragile]{Padrão Extend: \texttt{super()} + Extensão}
  Quando queremos \textbf{estender} (não substituir) o comportamento do pai:

  \begin{lstlisting}
class XTudoEspecial(XTudo):
    def __init__(self, preco, extras=None):
        super().__init__(preco, molho="da casa")
        self.extras = extras or ["bacon", "ovo"]

    def descrever(self):
        base = super().descrever()  # chama Lanche.descrever()
        return f"{base} + {', '.join(self.extras)}"
  \end{lstlisting}

  \begin{lstlisting}
xe = XTudoEspecial(32.0)
print(xe.descrever())
# X-Tudo: R$32.00 + bacon, ovo
  \end{lstlisting}
\end{frame}

\begin{frame}[fragile]{Armadilha: Esquecer \texttt{super().\_\_init\_\_()}}
  \begin{lstlisting}
class Coxinha(Lanche):
    def __init__(self, preco):
        # ESQUECEU super().__init__() !!!
        self.recheio = "frango"

c = Coxinha(8.0)
c.recheio       # "frango" --- OK
c.descrever()   # AttributeError: 'Coxinha' has no attribute 'nome'
  \end{lstlisting}

  \begin{alertblock}{Por quê?}
    Sem \texttt{super().\_\_init\_\_()}, os atributos \texttt{nome} e \texttt{preco} nunca são criados.
  \end{alertblock}

  \begin{exampleblock}{Correção}
    Sempre chame \texttt{super().\_\_init\_\_()} no início do \texttt{\_\_init\_\_} da subclasse.
  \end{exampleblock}
\end{frame}

% =============================================================================
\section{Polimorfismo e Duck Typing}
% =============================================================================

\begin{frame}[fragile]{Polimorfismo: Mesma Chamada, Comportamentos Diferentes}
  \textbf{Polimorfismo}: objetos de tipos diferentes respondem à mesma mensagem, cada um à sua maneira.

  \begin{lstlisting}
pedido = [XTudo(25.0), Dogao(18.0), Acai(15.0)]

for item in pedido:
    item.preparar()  # cada classe executa seu próprio preparar()
  \end{lstlisting}

  Saída:
  \begin{lstlisting}[numbers=none]
Fritando hambúrguer com molho especial...
Montando hot-dog grande...
Batendo açaí com granola, banana...
  \end{lstlisting}

  \begin{exampleblock}{Essência}
    O código \textbf{não precisa saber} o tipo exato --- basta que o objeto tenha o método.
  \end{exampleblock}
\end{frame}

\begin{frame}[fragile]{Duck Typing}
  \begin{quote}
    \emph{``If it walks like a duck and quacks like a duck, then it must be a duck.''}
  \end{quote}

  \bigskip

  \begin{lstlisting}
class Bebida:          # NAO herda de Lanche!
    def __init__(self, nome, preco):
        self.nome = nome
        self.preco = preco

    def preparar(self):
        print(f"Servindo {self.nome}...")

    def descrever(self):
        return f"{self.nome}: R${self.preco:.2f}"
  \end{lstlisting}
\end{frame}

\begin{frame}[fragile,shrink=10]{Duck Typing na Prática}
  \begin{lstlisting}
pedido = [XTudo(25.0), Dogao(18.0), Bebida("Guaraná", 6.0)]

for item in pedido:
    print(item.descrever())
    item.preparar()
  \end{lstlisting}

  Saída:
  \begin{lstlisting}[numbers=none]
X-Tudo: R$25.00
Fritando hambúrguer com molho especial...
Dogão: R$18.00
Montando hot-dog grande...
Guaraná: R$6.00
Servindo Guaraná...
  \end{lstlisting}

  \begin{alertblock}{Observação}
    \texttt{Bebida} não herda de \texttt{Lanche}, mas funciona no mesmo loop porque tem os mesmos métodos. Isso é \textbf{duck typing}.
  \end{alertblock}
\end{frame}

% =============================================================================
\section{Herança vs.\ Composição}
% =============================================================================

\begin{frame}{Herança vs.\ Composição}
  \begin{center}
    \begin{tabular}{lll}
      \toprule
      \textbf{Relacionamento} & \textbf{Mecanismo} & \textbf{Exemplo} \\
      \midrule
      ``É um'' (\emph{is-a})  & Herança    & XTudo \textbf{é um} Lanche \\
      ``Tem um'' (\emph{has-a}) & Composição & Pedido \textbf{tem} Lanches \\
      ``Tem um'' (\emph{has-a}) & Composição & Pedido \textbf{tem} FormaPagamento \\
      ``É um'' (\emph{is-a})  & Herança    & Pix \textbf{é uma} FormaPagamento \\
      \bottomrule
    \end{tabular}
  \end{center}
\end{frame}

\begin{frame}[fragile,shrink=10]{Composição no Código}
  \begin{lstlisting}
class FormaPagamento:
    def processar(self, valor):
        raise NotImplementedError

class Pix(FormaPagamento):
    def processar(self, valor):
        print(f"Pix de R${valor:.2f} confirmado!")

class Pedido:
    def __init__(self, pagamento):
        self.itens = []               # composição: tem lanches
        self.pagamento = pagamento    # composição: tem forma de pgto

    def adicionar(self, item):
        self.itens.append(item)

    def fechar(self):
        total = sum(i.preco for i in self.itens)
        self.pagamento.processar(total)
  \end{lstlisting}
\end{frame}

\begin{frame}[fragile]{Usando o Pedido}
  \begin{lstlisting}
pix = Pix()
pedido = Pedido(pix)
pedido.adicionar(XTudo(25.0))
pedido.adicionar(Bebida("Guaraná", 6.0))
pedido.fechar()
# Pix de R$31.00 confirmado!
  \end{lstlisting}

  \begin{exampleblock}{Gang of Four}
    \emph{``Favor composition over inheritance.''} \\
    Se você só quer \textbf{reusar} comportamento, sem relação conceitual ``é um'', use composição.
  \end{exampleblock}
\end{frame}

\begin{frame}{Quando Usar Cada Um?}
  \begin{center}
    \begin{tabular}{lcc}
      \toprule
      \textbf{Critério} & \textbf{Herança} & \textbf{Composição} \\
      \midrule
      Relação ``é um'' genuína & \checkmark & \\
      Apenas reusar código     &   & \checkmark \\
      Trocar comportamento em runtime & & \checkmark \\
      Hierarquia natural e estável & \checkmark & \\
      Flexibilidade máxima & & \checkmark \\
      \bottomrule
    \end{tabular}
  \end{center}

  \begin{alertblock}{Regra Prática}
    Comece com composição. Use herança apenas quando o ``é um'' for claro e estável.
  \end{alertblock}
\end{frame}

% =============================================================================
\section{Herança Múltipla e MRO}
% =============================================================================

\begin{frame}[fragile]{Herança Múltipla em Python}
  Python permite herdar de \textbf{mais de uma classe}:

  \begin{lstlisting}
class Promocional:
    def aplicar_desconto(self, pct):
        self.preco *= (1 - pct / 100)

class XTudoPromo(XTudo, Promocional):
    pass

xp = XTudoPromo(25.0)
xp.aplicar_desconto(10)
print(xp.preco)   # 22.5
  \end{lstlisting}
\end{frame}

\begin{frame}[fragile,shrink=10]{MRO --- Method Resolution Order}
  Quando há conflito, Python usa o \textbf{C3 Linearization} para decidir qual método chamar:

  \begin{lstlisting}
print(XTudoPromo.__mro__)
# (XTudoPromo, XTudo, Lanche, Promocional, object)
  \end{lstlisting}

  Python busca o método na ordem da MRO: da esquerda para a direita.

  \begin{alertblock}{Conselho}
    Herança múltipla pode causar confusão. Na prática:
    \begin{itemize}
      \item Use \textbf{mixins} simples (como \texttt{Promocional})
      \item Prefira \textbf{composição} para combinações complexas
      \item Evite hierarquias diamante
    \end{itemize}
  \end{alertblock}
\end{frame}

% =============================================================================
\section{Classes Abstratas}
% =============================================================================

\begin{frame}[fragile,shrink=10]{O Problema: Métodos ``Esquecidos''}
  \begin{lstlisting}
class Funcionario:
    def __init__(self, nome, salario_base):
        self.nome = nome
        self.salario_base = salario_base

    def calcular_salario(self):
        raise NotImplementedError("Subclasse deve implementar")
  \end{lstlisting}

  \begin{lstlisting}
class Atendente(Funcionario):
    pass   # esqueceu de implementar calcular_salario!

a = Atendente("Maria", 1500)
a.calcular_salario()  # NotImplementedError --- mas só descobre AQUI!
  \end{lstlisting}

  \begin{alertblock}{Problema}
    O erro só aparece quando \texttt{calcular\_salario()} é chamado --- pode ser tarde demais.
  \end{alertblock}
\end{frame}

\begin{frame}[fragile,shrink=15]{Solução: \texttt{ABC} e \texttt{@abstractmethod}}
  \begin{lstlisting}
from abc import ABC, abstractmethod

class Funcionario(ABC):
    def __init__(self, nome, salario_base):
        self.nome = nome
        self.salario_base = salario_base

    @abstractmethod
    def calcular_salario(self):
        """Cada tipo de funcionário calcula diferente."""
        pass
  \end{lstlisting}

  \begin{lstlisting}
# Tentar instanciar a classe abstrata:
f = Funcionario("João", 2000)
# TypeError: Can't instantiate abstract class Funcionario
#   with abstract method calcular_salario
  \end{lstlisting}

  \begin{exampleblock}{Vantagem}
    O erro ocorre na \textbf{instanciação}, não na chamada do método --- falha cedo!
  \end{exampleblock}
\end{frame}

\begin{frame}[fragile,shrink=10]{Subclasses Concretas}
  \begin{lstlisting}
class Atendente(Funcionario):
    def calcular_salario(self):
        return self.salario_base

class Cozinheiro(Funcionario):
    def __init__(self, nome, salario_base, adicional_insalubridade):
        super().__init__(nome, salario_base)
        self.adicional = adicional_insalubridade

    def calcular_salario(self):
        return self.salario_base + self.adicional

class Entregador(Funcionario):
    def __init__(self, nome, salario_base, entregas):
        super().__init__(nome, salario_base)
        self.entregas = entregas

    def calcular_salario(self):
        return self.salario_base + self.entregas * 5
  \end{lstlisting}
\end{frame}

\begin{frame}[fragile,shrink=20]{Loop Polimórfico para a Folha de Pagamento}
  \begin{lstlisting}
equipe = [
    Atendente("Maria", 1500),
    Cozinheiro("João", 1800, 300),
    Entregador("Carlos", 1200, 80),
    Atendente("Ana", 1500),
]

total_folha = 0
for func in equipe:
    salario = func.calcular_salario()
    total_folha += salario
    print(f"{func.nome:.<20} R${salario:>8.2f}")

print(f"{'TOTAL':.<20} R${total_folha:>8.2f}")
  \end{lstlisting}

  Saída:
  \begin{lstlisting}[numbers=none]
Maria............... R$ 1500.00
João................ R$ 2100.00
Carlos.............. R$ 1600.00
Ana................. R$ 1500.00
TOTAL............... R$ 6700.00
  \end{lstlisting}
\end{frame}

% #############################################################################
% PARTE 4: Tópicos Avançados e Boas Práticas
% #############################################################################

% ================================================================
\section{@dataclass --- Eliminando Boilerplate}
% ================================================================

\begin{frame}[fragile,shrink=25]{O Problema: Boilerplate Repetitivo}
Uma classe simples com \texttt{\_\_init\_\_}, \texttt{\_\_repr\_\_} e \texttt{\_\_eq\_\_}:

\begin{lstlisting}
class Lanche:
    def __init__(self, nome, preco, ingredientes):
        self.nome = nome
        self.preco = preco
        self.ingredientes = ingredientes

    def __repr__(self):
        return (f"Lanche(nome={self.nome!r}, "
                f"preco={self.preco!r}, "
                f"ingredientes={self.ingredientes!r})")

    def __eq__(self, other):
        if not isinstance(other, Lanche):
            return NotImplemented
        return (self.nome == other.nome
                and self.preco == other.preco
                and self.ingredientes == other.ingredientes)
\end{lstlisting}

\begin{alertblock}{Problema}
$\sim$20 linhas para guardar 3 campos. E se a classe tiver 8 campos?
\end{alertblock}
\end{frame}

\begin{frame}[fragile,shrink=3]{@dataclass: A Solução}
\begin{lstlisting}
from dataclasses import dataclass

@dataclass
class Lanche:
    nome: str
    preco: float
    ingredientes: list[str]
\end{lstlisting}

O decorador gera automaticamente:
\begin{itemize}
  \item \texttt{\_\_init\_\_} --- com todos os campos como parâmetros
  \item \texttt{\_\_repr\_\_} --- representação legível
  \item \texttt{\_\_eq\_\_} --- comparação campo a campo
\end{itemize}

\begin{exampleblock}{Resultado}
3 linhas substituem $\sim$20. O código expressa \textbf{o quê}, não \textbf{o como}.
\end{exampleblock}
\end{frame}

\begin{frame}[fragile]{Parâmetros Úteis do @dataclass}
\begin{lstlisting}
@dataclass(frozen=True)
class ItemCardapio:
    nome: str
    preco: float
    categoria: str
\end{lstlisting}

\begin{table}
\centering
\begin{tabular}{lll}
\toprule
\textbf{Parâmetro} & \textbf{Efeito} & \textbf{Quando usar} \\
\midrule
\texttt{frozen=True}  & Imutável (sem \texttt{setattr}) & Itens de cardápio, configs \\
\texttt{order=True}   & Gera \texttt{<, <=, >, >=}     & Ordenar por campos \\
\texttt{eq=True}       & Gera \texttt{\_\_eq\_\_} (padrão) & Quase sempre \\
\texttt{slots=True}   & Usa \texttt{\_\_slots\_\_} (3.10+) & Performance \\
\bottomrule
\end{tabular}
\end{table}
\end{frame}

\begin{frame}[fragile]{Defaults Mutáveis: field(default\_factory)}
\begin{lstlisting}
from dataclasses import dataclass, field

@dataclass
class Pedido:
    cliente: str
    itens: list[str] = field(default_factory=list)
    observacoes: dict[str, str] = field(default_factory=dict)
    numero: int = field(default=0, repr=False)
\end{lstlisting}

\begin{alertblock}{Nunca faça isso}
\texttt{itens: list[str] = []} --- o mesmo objeto é compartilhado entre instâncias!
\end{alertblock}

\texttt{field(default\_factory=list)} cria uma \textbf{nova lista} para cada instância.
\end{frame}

\begin{frame}[fragile,shrink=5]{@dataclass com Métodos Normais}
\begin{lstlisting}
@dataclass
class Pedido:
    cliente: str
    itens: list[str] = field(default_factory=list)

    def adicionar(self, item: str):
        self.itens.append(item)

    def total(self, cardapio: dict[str, float]) -> float:
        return sum(cardapio[i] for i in self.itens)

    def resumo(self) -> str:
        return f"{self.cliente}: {len(self.itens)} itens"
\end{lstlisting}

\begin{exampleblock}{Ponto-chave}
\texttt{@dataclass} gera o boilerplate; você adiciona a lógica de negócio.
\end{exampleblock}
\end{frame}

\begin{frame}{Quando Usar / Quando NÃO Usar @dataclass}
\begin{columns}[T]
\column{0.48\textwidth}
\textbf{Use quando:}
\begin{itemize}
  \item Classe é principalmente ``dados com métodos''
  \item Precisa de \texttt{\_\_eq\_\_} e \texttt{\_\_repr\_\_}
  \item Muitos campos (reduz boilerplate)
  \item Quer imutabilidade (\texttt{frozen})
\end{itemize}

\column{0.48\textwidth}
\textbf{NÃO use quando:}
\begin{itemize}
  \item \texttt{\_\_init\_\_} tem lógica complexa (validação pesada)
  \item Herança profunda ($>$2 níveis)
  \item Classe é majoritariamente comportamento
  \item Compatibilidade com Python $<$3.7
\end{itemize}
\end{columns}
\end{frame}

% ================================================================
\section{@classmethod e @staticmethod}
% ================================================================

\begin{frame}[shrink=5]{Três Tipos de Método}
\begin{table}
\centering
\begin{tabular}{llll}
\toprule
\textbf{Tipo} & \textbf{Decorador} & \textbf{1\textsuperscript{o} Parâmetro} & \textbf{Acessa} \\
\midrule
Instância   & (nenhum)          & \texttt{self} & instância + classe \\
Classe      & \texttt{@classmethod}  & \texttt{cls}  & apenas classe \\
Estático    & \texttt{@staticmethod} & (nenhum)      & nada da classe \\
\bottomrule
\end{tabular}
\end{table}

\begin{exampleblock}{Regra prática}
\textbf{Precisa de \texttt{self}?} $\to$ método de instância.
\textbf{Precisa de \texttt{cls}?} $\to$ classmethod.
\textbf{Não precisa de nenhum?} $\to$ staticmethod.
\end{exampleblock}
\end{frame}

\begin{frame}[fragile,shrink=10]{@classmethod como Construtor Alternativo}
\begin{lstlisting}
@dataclass
class Lanche:
    nome: str
    preco: float
    categoria: str

    @classmethod
    def from_csv_line(cls, linha: str) -> "Lanche":
        nome, preco, cat = linha.strip().split(",")
        return cls(nome, float(preco), cat)

    @classmethod
    def from_dict(cls, d: dict) -> "Lanche":
        return cls(**d)
\end{lstlisting}

\begin{lstlisting}
x = Lanche.from_csv_line("X-Tudo,25.90,hamburguer")
y = Lanche.from_dict({"nome": "Suco", "preco": 8.0,
                       "categoria": "bebida"})
\end{lstlisting}
\end{frame}

\begin{frame}[fragile,shrink=3]{@staticmethod para Utilitários}
\begin{lstlisting}
class Lanche:
    # ...

    @staticmethod
    def calcular_taxa_servico(valor: float,
                               percentual: float = 0.10) -> float:
        """Taxa de serviço sobre o valor."""
        return round(valor * percentual, 2)
\end{lstlisting}

\begin{lstlisting}
taxa = Lanche.calcular_taxa_servico(50.00)  # 5.0
taxa = Lanche.calcular_taxa_servico(50.00, 0.15)  # 7.5
\end{lstlisting}

\begin{exampleblock}{Quando usar}
Lógica relacionada à classe, mas que não precisa de \texttt{self} nem de \texttt{cls}. Alternativa: função no módulo.
\end{exampleblock}
\end{frame}

\begin{frame}[shrink=25]{Guia de Decisão}
\begin{center}
\begin{tikzpicture}[
  node distance=1.2cm,
  every node/.style={font=\small},
  decision/.style={diamond, draw, fill=blue!10, text width=3.2cm,
    align=center, inner sep=1pt, aspect=2},
  result/.style={rectangle, draw, fill=green!10, rounded corners,
    text width=3cm, align=center, minimum height=0.8cm},
  arr/.style={->, thick}
]
\node[decision] (q1) {Precisa acessar\\a instância?};
\node[result, right=2.5cm of q1] (r1) {Método de\\instância (\texttt{self})};
\node[decision, below=1.2cm of q1] (q2) {Precisa acessar\\a classe?};
\node[result, right=2.5cm of q2] (r2) {\texttt{@classmethod}\\(\texttt{cls})};
\node[result, below=1.2cm of q2] (r3) {\texttt{@staticmethod}};

\draw[arr] (q1) -- node[above]{Sim} (r1);
\draw[arr] (q1) -- node[left]{Não} (q2);
\draw[arr] (q2) -- node[above]{Sim} (r2);
\draw[arr] (q2) -- node[left]{Não} (r3);
\end{tikzpicture}
\end{center}
\end{frame}

% ================================================================
\section{Princípios SOLID (Introdução)}
% ================================================================

\begin{frame}{SOLID --- Visão Geral}
\begin{table}
\centering
\small
\begin{tabular}{clp{6cm}}
\toprule
\textbf{Letra} & \textbf{Princípio} & \textbf{Ideia Central} \\
\midrule
\textbf{S} & Single Responsibility  & Uma classe, uma razão para mudar \\
\textbf{O} & Open/Closed            & Aberta p/ extensão, fechada p/ modificação \\
\textbf{L} & Liskov Substitution    & Subtipos substituem o pai sem surpresas \\
\textbf{I} & Interface Segregation  & Interfaces pequenas e coesas \\
\textbf{D} & Dependency Inversion   & Dependa de abstrações, não de concretos \\
\bottomrule
\end{tabular}
\end{table}

Veremos \textbf{S}, \textbf{O} e \textbf{L} com exemplos do Podrão.
\end{frame}

\begin{frame}[fragile]{SRP --- O Problema}
\begin{lstlisting}
class PedidoRuim:
    def __init__(self, cliente, itens):
        self.cliente = cliente
        self.itens = itens

    def calcular_total(self): ...       # lógica de negócio
    def gerar_relatorio_html(self): ... # apresentação
    def salvar_em_arquivo(self): ...    # persistência
    def enviar_email(self): ...         # comunicação
\end{lstlisting}

\begin{alertblock}{4 razões para mudar}
Mudança no cálculo, no formato HTML, no storage, ou no e-mail --- todas alteram a mesma classe.
\end{alertblock}
\end{frame}

\begin{frame}[fragile]{SRP --- A Solução: Separar Responsabilidades}
\begin{lstlisting}
class Pedido:
    def calcular_total(self): ...

class RelatorioHTML:
    def gerar(self, pedido: Pedido): ...

class Armazenamento:
    def salvar(self, pedido: Pedido): ...

class Enviador:
    def enviar_email(self, pedido: Pedido, dest: str): ...
\end{lstlisting}

\begin{exampleblock}{Insight}
Cada classe tem \textbf{uma única razão para mudar}. Testar cada uma é simples e independente.
\end{exampleblock}
\end{frame}

\begin{frame}[fragile]{OCP --- O Problema}
\begin{lstlisting}
class CalculadoraDesconto:
    def calcular(self, pedido, tipo):
        if tipo == "percentual":
            return pedido.total() * 0.10
        elif tipo == "fixo":
            return 5.00
        elif tipo == "fidelidade":
            return pedido.total() * 0.15
        # ... cada novo tipo exige elif
\end{lstlisting}

\begin{alertblock}{Violação OCP}
Para adicionar ``combo'', precisa \textbf{modificar} \texttt{CalculadoraDesconto}. A classe não está fechada para modificação.
\end{alertblock}
\end{frame}

\begin{frame}[fragile,shrink=20]{OCP --- A Solução: Polimorfismo}
\begin{lstlisting}
from abc import ABC, abstractmethod

class Desconto(ABC):
    @abstractmethod
    def calcular(self, total: float) -> float: ...

class DescontoPercentual(Desconto):
    def __init__(self, taxa: float):
        self.taxa = taxa
    def calcular(self, total: float) -> float:
        return total * self.taxa

class DescontoFixo(Desconto):
    def __init__(self, valor: float):
        self.valor = valor
    def calcular(self, total: float) -> float:
        return min(self.valor, total)
\end{lstlisting}

\begin{exampleblock}{Insight}
Novo desconto = nova classe. \texttt{Pedido} não muda. \textbf{Extensão sem modificação.}
\end{exampleblock}
\end{frame}

\begin{frame}[fragile,shrink=15]{LSP --- O Problema Clássico}
\begin{lstlisting}
class Retangulo:
    def __init__(self, largura, altura):
        self._largura = largura
        self._altura = altura

    def set_largura(self, v):
        self._largura = v

    def area(self):
        return self._largura * self._altura

class Quadrado(Retangulo):
    def set_largura(self, v):
        self._largura = v
        self._altura = v   # surpresa!
\end{lstlisting}

\begin{alertblock}{Violação LSP}
Código que espera \texttt{Retangulo} e chama \texttt{set\_largura(5)} seguido de \texttt{area()} obtém resultado inesperado com \texttt{Quadrado}.
\end{alertblock}
\end{frame}

% ================================================================
\section{Padrões de Projeto}
% ================================================================

\begin{frame}[fragile,shrink=15]{Strategy Pattern --- Composição de Comportamento}
\begin{lstlisting}
class SemDesconto(Desconto):
    def calcular(self, total: float) -> float:
        return 0.0

@dataclass
class Pedido:
    cliente: str
    itens: list[str] = field(default_factory=list)
    desconto: Desconto = field(
        default_factory=SemDesconto)

    def total(self, cardapio: dict) -> float:
        bruto = sum(cardapio[i] for i in self.itens)
        return bruto - self.desconto.calcular(bruto)
\end{lstlisting}

\begin{exampleblock}{Composição}
\texttt{Pedido} recebe a estratégia; não sabe qual é. Novo desconto = nova classe, zero mudança em \texttt{Pedido}.
\end{exampleblock}
\end{frame}

\begin{frame}[fragile]{Strategy --- Uso}
\begin{lstlisting}
cardapio = {"X-Tudo": 25.90, "Suco": 8.00}

p1 = Pedido("Ana", desconto=DescontoPercentual(0.10))
p1.adicionar("X-Tudo")
p1.adicionar("Suco")
print(p1.total(cardapio))  # 30.51

p2 = Pedido("Bruno", desconto=DescontoFixo(5.00))
p2.adicionar("X-Tudo")
print(p2.total(cardapio))  # 20.90
\end{lstlisting}

A estratégia pode ser trocada em \textbf{tempo de execução}:
\begin{lstlisting}
p1.desconto = DescontoFixo(10.00)
\end{lstlisting}
\end{frame}

\begin{frame}[fragile,shrink=5]{Factory Method --- Registry}
\begin{lstlisting}
class LancheFactory:
    _registry: dict[str, type] = {}

    @classmethod
    def registrar(cls, nome: str, classe: type):
        cls._registry[nome] = classe

    @classmethod
    def criar(cls, nome: str, **kwargs):
        if nome not in cls._registry:
            raise ValueError(f"Tipo desconhecido: {nome}")
        return cls._registry[nome](**kwargs)

# Registro
LancheFactory.registrar("xtudo", XTudo)
LancheFactory.registrar("xsalada", XSalada)
LancheFactory.registrar("hamburguer", Hamburguer)
\end{lstlisting}
\end{frame}

\begin{frame}[fragile,shrink=5]{Factory Method --- Uso e Extensão}
\begin{lstlisting}
# Criação via string (ex: lido de arquivo/API)
lanche = LancheFactory.criar("xtudo", preco=25.90)

# Extensão em runtime --- sem tocar na Factory
class XBacon(Lanche):
    def __init__(self, preco=28.00):
        super().__init__("X-Bacon", preco, ["bacon"])

LancheFactory.registrar("xbacon", XBacon)
lanche2 = LancheFactory.criar("xbacon")
\end{lstlisting}

\begin{exampleblock}{Vantagens}
\begin{itemize}
  \item Desacopla criação de uso
  \item Extensível em runtime (\texttt{registrar})
  \item String $\to$ objeto (útil para configs, APIs, CLIs)
\end{itemize}
\end{exampleblock}
\end{frame}

% ================================================================
\section{Resumo do Módulo}
% ================================================================

\begin{frame}{Mapa de Conceitos: POO Completo}
\begin{table}
\centering
\scriptsize
\begin{tabular}{lll}
\toprule
\textbf{Conceito} & \textbf{Seção} & \textbf{No Podrão} \\
\midrule
Classe e objeto          & Classes e Objetos & \texttt{Lanche}, \texttt{Pedido} \\
Atributos e métodos      & Classes e Objetos & \texttt{nome}, \texttt{preco()}, \texttt{adicionar()} \\
Encapsulamento           & Encapsulamento & \texttt{\_preco}, \texttt{@property} \\
Dunders                  & Métodos Especiais & \texttt{\_\_str\_\_}, \texttt{\_\_eq\_\_}, \texttt{\_\_add\_\_} \\
Herança                  & Herança Simples & \texttt{XTudo(Lanche)}, \texttt{Dogao(Lanche)} \\
Classe abstrata          & Classes Abstratas & \texttt{Funcionario(ABC)}, \texttt{Desconto(ABC)} \\
Polimorfismo             & Polimorfismo & \texttt{preparar()} em cada subtipo \\
Composição               & Herança vs Composição & \texttt{Pedido} tem lista de \texttt{Lanche} \\
\texttt{@dataclass}      & @dataclass & \texttt{Pedido}, \texttt{ItemCardapio} \\
\texttt{@classmethod}    & @classmethod & \texttt{Lanche.from\_csv\_line()} \\
Strategy                 & Padrões de Projeto & \texttt{Desconto} $\to$ \texttt{DescontoPercentual} \\
Factory                  & Padrões de Projeto & \texttt{LancheFactory.criar()} \\
SOLID (S, O, L)          & Princípios SOLID & SRP, OCP, LSP \\
\bottomrule
\end{tabular}
\end{table}
\end{frame}

\begin{frame}{Os Três Mantras da POO}
\begin{enumerate}
  \item \textbf{Objetos agrupam dados e comportamento}

  Não use listas paralelas. Crie uma classe que una o dado ao que se faz com ele.

  \item \textbf{Prefira composição a herança}

  Herança cria acoplamento forte. Composição permite trocar comportamento em runtime (Strategy).

  \item \textbf{Cada classe, uma responsabilidade + extensão sem modificação}

  SRP + OCP. Classes pequenas, coesas, extensíveis via polimorfismo.
\end{enumerate}

\begin{exampleblock}{Se lembrar apenas disso\ldots}
Estes três princípios guiam a maior parte das decisões de design em POO.
\end{exampleblock}
\end{frame}

\end{document}
